\section{Future work}
\label{sec:future}

Future work on \DEIC{} is distributed along two axes: technical extensions (methodological, requiring no theoretical rework) and fundamental continuations (the elaboration of the theoretical gaps from §\ref{sec:limitations}, each of which becomes a separate paper in the program).

\subsection{Technical extensions}

\paragraph{Pre-registered replication on a new sample.} A mandatory condition prior to wide publication. The central predictions: A-moderation of the CA$\to$ID$\to$\EAF{} channel in the narrow specification; collapse of the ID path in the unified SEM upon inclusion of P; the two-axis distinction of the primary and secondary patterns; the compensatory function of A in the hi-CA$\times$hi-A group; the positive partial correlation of P and \Ecog{} after partialling. Replication must reproduce this combination specifically, including the re-composition of variance in the unified model.

\paragraph{Item wave 2: instrument rebuild.} Priority directions: (1) a dedicated A-scale, separated by design between witnessing and trauma hypervigilance, with content validation across all four linguistic traditions (satipa\d{t}\d{t}h\=ana, s\=ak\d{s}\=i, rigpa, self-remembering); (2) M-redesign separating instrumental and reactive masking; (3) \EAF{}/\Ecog{} separated at the item level rather than post-hoc; (4) one item $=$ one factor in scoring (without multi-scale weights); (5) P scale extended with an explicit moral signature (items on ``it is allowed for me'') separately from architectural items on resonance switching; (6) ACE items retain a three-option format (``yes'' / ``no'' / ``don't remember''), but ``don't remember'' is coded as an explicit \emph{positive} signal (family silencing or dissociative encoding gap), not as missing; the count of such responses across all 10 items forms a pre-registered ACE-dissociation subscale with stated criteria of McDonald's $\omega \ge 0.7$ and $r(\text{ACE-dis}, \text{ID}_{\mathrm{lat}}) \ge 0.30$. Post-hoc re-fit on wave 1 (§\ref{subsec:ace_extension}) shows strengthening of the latent association $r(\text{CA}, \EAF{})$ from $-0.129$ to $-0.205$ and the emergence of a statistically nonzero indirect path through P---wave 2 must pre-register this logic and confirm it as an established fact.

\paragraph{Longitudinal extension} (6--12--24 months, in-therapy subgroup). Tests the stability of A, the trajectory of integration, and the test--retest reliability of all six factors. A sample with explicit stratification by therapy-status (in-treatment, post-treatment, never-treated) will permit a distinction between stable and transient A.

\paragraph{Multicultural validation} (EN, ES, CH, AR) with invariance testing (configural $\to$ metric $\to$ scalar). A test of the generalizability of the architecture beyond a Russian-language sample.

\paragraph{fMRI convergence.} DMN $\times$ attention-networks co-activation as a neural correlate of A. Biological grounding of the A construct via existing neuroimaging paradigms (resting-state functional connectivity, mindfulness-induced DMN deactivation).

\paragraph{Informant report on a subsample.} Round-robin validity is especially important for A and P, where self-report has structural limitations.

\paragraph{Behavioral convergents for A.} Dual-task paradigms, interoception accuracy tasks (heartbeat detection), mind-wandering probes---an exit from pure self-report.

\subsection{Fundamental continuations}

These are papers of the program developing the theoretical gaps of §\ref{sec:limitations}.

\paragraph{``\DEIC{} dynamics: temporal architecture of empathy blocking.''} A three-timescale framework (fast/meso/slow), connection with active inference and predictive coding. Empirical correlates via micro-longitudinal (daily diary) and neurocognitive (ERP) paradigms.

\paragraph{``ID and P as directional asymmetries of empathic prediction.''} Separation of bottom-up ID and top-down P within the predictive-coding frame. Interoceptive accuracy tasks as the differentiator: ID patients should show a deficit in interoceptive accuracy; P patients---preserved interoception with reduced moral salience of others' distress.

\paragraph{``What does attention do: candidate mechanisms for the A-moderator effect.''} A three-channel theoretical paper: temporal resolution, automatic-binding decomposition, metacognitive working memory. Empirical probes via phasic attention tasks and dual-task paradigms.

\paragraph{``Individual differences in A-trainability: who benefits from witnessing-based therapy?''} Pre-treatment predictors of the effect of A-based interventions. Candidates: working memory capacity, interoceptive accuracy, language-mediation, safety-context stability. A clinical sample with an RCT arm.

\paragraph{``Developmental branching in the $2\times 2$ psychopathic topology.''} Longitudinal child sample with early CA/temperament markers; a test of the Karpman hypothesis of two routes into the primary vs.\ secondary cell.

\paragraph{``\DEIC{} and the reflexivity of measurement: observer-A as part of the instrument.''} A methodological paper on the claim that the observer's A is a condition of possibility for seeing A in the observed. Epistemologically: a structural, not mystical, analogy with the observer problem in quantum mechanics. Operational consequences: supervision requirements for clinicians, A-assessment of the researcher as part of methodological reporting.

\paragraph{``Free will as a gradient of attention: an operational reading.''} A separate philosophical-psychological paper developing §\ref{subsec:free_will}. Proposes a reformulation of the problem of free will as a problem of measurable A-infrastructure, with testable predictions: the correlation between the felt sense of free choice and momentary A (rather than trait self-efficacy), increases in self-determination markers following A-interventions, structural functional unfreedom in low-A populations. Unpacks the ethical consequences (a recalibration of responsibility, forgiveness as epistemological correctness) without erasing the distinction between low-A automatism and A-supported P choice.

\paragraph{``Pharmacology as \DEIC{} modification: a stratification-based prediction map.''} A paper formalizing the table ``class of substance---modification of component D/P/A---profile-based prediction---proposed clinical test,'' developed in §\ref{subsec:pharmacology_nosology}. The core: differential predictions: psilocybin with a large effect on D-dominant, small/reverse on P-dominant; MDMA unique as E-ceiling $+$ D-dissolve $+$ A-preserve; ketamine as ``dissociation for the dissolution of dissociation'' with responders $=$ hi-D baseline; opioid dependence as a safety-deficit-D signature. Proposed empirical program: retrospective re-analysis of existing psilocybin/MDMA RCTs with DEIC profile at baseline as a moderator of response.

\paragraph{``The structural under-sizing of modern mental healthcare: what early-formation D requires.''} A clinical-institutional paper developing §\ref{subsec:undersizing}. Core: an argument for a structurally different scale of intervention for personality-formative early D, with a review of methods of adequate scale (long-term intensive psychoanalysis, psychedelic-assisted protocols with extended integration, therapeutic communities, multi-year somatic work) and a map of ``what is currently unavailable.'' Not a manifesto; an empirical program of proof-of-concept: DEIC stratification at baseline as a predictor of response to symptom-level vs.\ structural-level interventions.

\paragraph{``\DEIC{} and population-level predictions: attention economy, politics, collective D.''} A paper developing §\ref{subsec:political}. Predictions: (a) intensity of use of engagement-optimized platforms reduces A in 6--12-month longitudinal windows; (b) adolescent cohorts with early algorithm exposure have a measurably different DEIC topology at the same ACE load; (c) electoral support of P-style demagoguery is predicted by the general A level of the audience; (d) DEIC profiles of ex-cult and ex-radicalized individuals are characterized by low A with preserved suppressed D. Data for (a)--(d) require cross-cohort and cross-cultural designs.

\paragraph{``Evolutionary mismatch and the biological necessity of integration.''} A programmatic paper developing §\ref{subsec:evolutionary}. Thesis: the contemporary adult is the first in the history of the species for whom integration is a biological requirement rather than a cultural luxury. Empirical program: population-level trajectories of A indicators (distress tolerance, boredom capacity, sustained attention) in relation to urban engagement density; incidence of D-dominant disorders as a function of the gap between ACE load and the availability of integration infrastructure; 20--40-year cohort comparisons of cohorts with and without structural access to integration systems.

\paragraph{``Nosological reorganization: DEIC profile as a trans-diagnostic stratifier.''} An empirical paper testing the capacity of the DEIC profile at baseline to capture a larger share of the variance of clinical response and trajectory than the DSM label, in a large-scale retrospective re-analysis. Compatibility with and parallel to HiTOP and RDoC; falsification: if the DEIC profile does not yield incremental gain over diagnosis label---the framework is restricted to narrower applications, the nosological claim is withdrawn.

\subsection{Applied directions}

\paragraph{Clinical RCTs.} Head-to-head tests: non-dual mindfulness vs.\ MBSR with ID-dominant clients; an A-training arm vs.\ 12-step in substance-use with a blocking profile; sequenced A-work $\to$ CBT vs.\ CBT-first in treatment-resistant depression.

\paragraph{DEIC as a solution to the under-diagnosis problem (collaborative-filtering approach).} A direct extension of §\ref{subsec:dx_count_bias}. The observation that the instrument systematically captures precisely those the self-report diagnostic misses (hi-P---through ROC AUC $< 0.5$ on mood/anxiety/neurodev; hi-ID---through help-seeking block; early-trauma hi-CA---through ``don't know'' as a dissociative signal) is inverted into a positive design. Classical nosology predicts diagnosis from reported symptoms; DEIC coordinates are stable in architecturally similar groups, so $k$-NN / latent-class similarity on the 6-factor profile permits prediction of a \emph{probabilistic clinical landscape} for a new respondent by proximity to a subsample with known diagnoses---the standard formulation of collaborative filtering but with a clinical target variable. Design for validation: a wave with recruitment via primary care and educational institutions (outside of the meme-channel, to break the selection bias), DEIC at intake, structured clinical interview (SCID/MINI) blind to the DEIC profile as gold standard; metrics: sensitivity/specificity of the DEIC screen against clinical labels, incremental AUC over symptom screening (PHQ, GAD), calibration curves for ``architecturally conditioned'' prognosis. The public meaning is subtler than it sounds: the goal is not automatic labeling of people, but removing the bottleneck on recognition where the very mechanism of pathology blocks recognition. The limits are evident (risk of false-positive labeling, the ethics of screening without consent, the necessity of a human-in-the-loop clinician) and must be the subject of a separate applied framework.

\paragraph{Public communication.} A separate frame ``trauma is not a verdict'' is being developed outside the academic layer and requires its own empirical validation on a non-academic audience. Theses: the integrated survivor as an empirically grounded horizon of recovery; the distinction between wellness-industry toxicity and contemplative substance; A-work as a path, not a promise of happiness.

\paragraph{Political and organizational communication.} P-style demagoguery and low-A audiences---a testable hypothesis via cross-sectional audience studies. Organizational psychology: the applicability of the conditional-architectural frame in the selection and supervision of roles requiring emotional regulation (medicine, helping professions, security services).

\paragraph{Education.} A-oriented school programs as primary prevention for CA-high cohorts. A connection with existing SEL (social-emotional learning) initiatives; empirical testing of whether SEL programs work on the mechanism that \DEIC{} posits.

\paragraph{Attention-protective institutional infrastructure.} An applied continuation of §\ref{subsec:political}. Design and evaluation of interventions at the level of media product design (slow-feed, pause interfaces, edit mode), at the level of school programs (contemplative education as a component of the general curriculum), at the level of regulation (engagement optimization as a public health issue). Testable: whether such changes yield a measurable increase in population-level A, other things equal.

\paragraph{Witnessing work on collective D: transitional justice as a DEIC operation.} An applied paper on how DEIC coordinates can give a quantitative measure to what transitional justice and memorial practices do intuitively: collective D-work through the creation of a public space for affective integration of historical reality. Pilot data: post-totalitarian societies with accessible longitudinal social-psychological data.

\subsection{Meta-program}

\paragraph{Methodological paper of the program.} A formal exposition of the anti-Wilber disciplinary rules (§\ref{subsec:discipline}), the reflexivity clause, the observer-A problem. Placement: paper \#18 in the program, issued after the empirical core.

\paragraph{Open documentation of the program.} \texttt{PROGRAM\_STRUCTURE.md} in the repository contains the full list of 18 planned papers with their falsifiers and expected deliverables. Each closes either on a pre-registered replication or on an explicit update of the program under non-replication. This is the form in which we accept commitments: a distributed-over-time network of testable claims, not a completed theory.

The present paper---narrow in tone, modest in its claims---defends the possibility of unfolding an ambitious program. \DEIC{} as a framework stands on the principle that each of its future works may be refuted at its own point, and if a significant share is refuted, the framework as a whole is subject to revision. Otherwise we are not working in psychometrics but in rhetoric.
